\chapter{Planteamiento del Problema} \label{sec:Planteamiento}
La adopción de Kubernetes como plataforma de orquestación de contenedores ha crecido de manera acelerada en empresas de distintos tamaños, convirtiéndose en un estándar de facto para el despliegue de aplicaciones modernas. Sin embargo, el éxito de esta tecnología depende en gran medida del plano de red. De este plano dependen la conectividad, el rendimiento, la seguridad y la eficiencia de los servicios. Kubernetes delega esta responsabilidad en los complementos de red (CNI), los cuales ofrecen múltiples alternativas con diferentes capacidades y costos operativos.

El reto para las organizaciones es que no existe un modelo práctico y claramente estructurado que permita seleccionar, implementar y mantener estas soluciones de manera objetiva. Las decisiones suelen basarse en pruebas limitadas, en documentación dispersa o en experiencias anecdóticas, lo que conduce a elecciones que no reflejan las condiciones reales de operación. En consecuencia, se generan configuraciones inseguras, políticas de red incompletas y sobredimensionamiento de recursos, problemas que impactan directamente en el rendimiento de los clústeres y en los costos de operación.

A nivel operativo, la falta de criterios de comparación dificulta la selección adecuada del CNI; la escasez de competencias especializadas retrasa implementaciones y aumenta la dependencia de terceros; y la ausencia de observabilidad sobre el tráfico de red provoca que los fallos se identifiquen únicamente cuando afectan al usuario final. Estos vacíos no solo degradan la calidad del servicio, sino que incrementan la exposición a riesgos de seguridad y generan gastos innecesarios en infraestructura y consumo energético.

La problemática se intensifica porque, aunque Kubernetes provee mecanismos como las NetworkPolicies o el autoescalado de recursos, su aplicación práctica requiere conocimientos avanzados y pruebas sistemáticas que la mayoría de equipos no realiza. En entornos donde los tiempos de entrega y la disponibilidad son críticos, operar bajo supuestos en lugar de evidencias convierte la red en un punto débil.

En este contexto, surge la necesidad de investigar cómo comparar de forma objetiva distintos CNIs y cómo automatizar la validación de políticas de red, de modo que las organizaciones puedan garantizar seguridad, rendimiento y eficiencia sin incrementar la complejidad de la operación. Resolver esta brecha es clave para que Kubernetes cumpla su promesa de escalabilidad y confiabilidad, y para que las empresas reduzcan riesgos y costos asociados a la infraestructura tecnológica.